
\documentclass{article}
\usepackage{colortbl}
\usepackage{makecell}
\usepackage{multirow}
\usepackage{supertabular}

\begin{document}

\newcounter{utterance}

\centering \large Interaction Transcript for game `hot\_air\_balloon', experiment `air\_balloon\_survival\_hu\_complexity\_hard', episode 1 with mistral{-}large{-}3{-}azure.
\vspace{24pt}

{ \footnotesize  \setcounter{utterance}{1}
\setlength{\tabcolsep}{0pt}
\begin{supertabular}{c@{$\;$}|p{.15\linewidth}@{}p{.15\linewidth}p{.15\linewidth}p{.15\linewidth}p{.15\linewidth}p{.15\linewidth}}
    \# & $\;$A & \multicolumn{4}{c}{Game Master} & $\;\:$B\\
    \hline

    \theutterance \stepcounter{utterance}  
    & & \multicolumn{4}{p{0.6\linewidth}}{
        \cellcolor[rgb]{0.9,0.9,0.9}{
            \makecell[{{p{\linewidth}}}]{
                \texttt{\tiny{[P1$\langle$GM]}}
                \texttt{Egy együttműködésen alapuló tárgyalós játékban veszel részt.} \\
\\ 
\texttt{Egy másik résztvevővel közösen kell megegyeznetek abban, hogy mely elemeket tartjátok meg. Mindkettőtöknek megvan a saját nézete arról, hogy számára mennyire fontosak az egyes elemek (fontosság). Nem tudod, a másik résztvevő hogyan értékeli az elemeket. Emellett minden elemhez tartozik egy erőfeszítés{-}érték is.} \\
\texttt{Csak akkor egyezhettek meg egy elemkészletben, ha a kiválasztott elemek összesített erőfeszítése nem haladja meg a közös korlátot:} \\
\\ 
\texttt{KORLÁT = 4497} \\
\\ 
\texttt{Az egyes elemekhez tartozó erőfeszítés{-}értékek:} \\
\\ 
\texttt{Elemek erőfeszítése = \{"C08": 828, "C55": 250, "A43": 412, "A67": 75, "B42": 275, "A62": 565, "B19": 889, "C43": 73, "C73": 747, "C10": 77, "B36": 23, "A18": 651, "C61": 11, "C67": 298, "C42": 769, "C70": 812, "B40": 368, "C39": 506, "A21": 481, "C16": 884\}} \\
\\ 
\texttt{Itt van, hogy számodra mennyire fontosak az egyes elemek:} \\
\\ 
\texttt{Elempreferencia értékek = \{"C08": 916, "C55": 338, "A43": 500, "A67": 163, "B42": 363, "A62": 653, "B19": 977, "C43": 161, "C73": 835, "C10": 165, "B36": 111, "A18": 739, "C61": 99, "C67": 386, "C42": 857, "C70": 900, "B40": 456, "C39": 594, "A21": 569, "C16": 972\}} \\
\\ 
\texttt{Cél:} \\
\\ 
\texttt{A célod, hogy olyan közös elemkészletet alkudj ki, amely számodra a lehető legelőnyösebb (azaz a számodra maximális összfontosságot eredményezi), miközben a készlet összerőfeszítése nem haladja meg a KORLÁT{-}ot. Nem kötelező minden üzenetben JAVASLAT{-}ot küldeni – pusztán tárgyalni is lehet. Minden taktika megengedett!} \\
\\ 
\texttt{Interakciós protokoll:} \\
\\ 
\texttt{Üzeneteidben kizárólag az alábbi struktúrált formátumokat használhatod:} \\
\\ 
\texttt{JAVASLAT: \{'A', 'B', 'C', …\}} \\
\texttt{Pontosan ezeknek az elemeknek a megtartását javasold.} \\
\texttt{ELUTASÍTÁS: \{'A', 'B', 'C', …\}} \\
\texttt{Kifejezetten utasítsd el a másik fél javaslatát.} \\
\texttt{ÉRVELÉS: \{'...'\}} \\
\texttt{Védd meg a saját javaslatodat, vagy érvelj a másik játékos javaslata ellen.} \\
\texttt{ELFOGADÁS: \{'A', 'B', 'C', …\}} \\
\texttt{Fogadd el a másik fél javaslatát, ezzel véget ér a játék.} \\
\\ 
\\ 
\texttt{Szabályok:} \\
\\ 
\texttt{Csak olyan javaslatot fogadhatsz el ELFOGADÁS formátumban, amelyet a másik fél JAVASLAT segítségével tett.} \\
\texttt{Csak olyan javaslatot utasíthatsz el ELUTASÍTÁS formátumban, amelyet a másik fél JAVASLAT{-}tal tett.} \\
\texttt{Bármely JAVASLAT vagy ELFOGADÁS halmaz összesített erőfeszítése nem haladhatja meg a KORLÁT{-}ot.} \\
\texttt{Ne áruld el a saját fontossági pontjaidat.} \\
\texttt{A strukturált formátumban szereplő címkét kettőspont és szóköz kell, hogy kövesse. Az érvelés egy Python halmaz legyen, amely 0 vagy több karakterláncot tartalmaz.} \\
\texttt{A formátumnak tehát így kell kinéznie: "CÍMKE: \{...\}"} \\
\texttt{Szigorúan tartsd be az interakciós protokollt, és NE írj semmi mást a megadott formátumokon kívül.} \\
\texttt{A játék akkor ér véget, ha az egyik fél ELFOGADÁS{-}t ad egy másik fél által JAVASLAT{-}tal tett javaslatra.} \\
\texttt{A válaszodnak olyan tartalmat is kell tartalmaznia, amely továbbadható a másik játékosnak – nem lehet üres.} \\
\texttt{Csak azok a javaslatok aktívak, amelyeket a másik játékos JAVASLAT segítségével tett, és amelyeket még nem utasítottál el ELUTASÍTÁS segítségével.} \\
\\ 
            }
        }
    }
    & \\ \\

    \theutterance \stepcounter{utterance}  
    & \multicolumn{4}{p{0.6\linewidth}}{
        \cellcolor[rgb]{0.9,0.9,0.9}{
            \makecell[{{p{\linewidth}}}]{
                \texttt{\tiny{[P1$\rangle$GM]}}
                \texttt{JAVASLAT: \{'C08', 'C55', 'A43', 'A67', 'B42', 'A62', 'C43', 'C10', 'B36', 'C61', 'C67', 'B40', 'A21'\}} \\
            }
        }
    }
    & & \\ \\

    \theutterance \stepcounter{utterance}  
    & & \multicolumn{4}{p{0.6\linewidth}}{
        \cellcolor[rgb]{0.9,0.9,0.9}{
            \makecell[{{p{\linewidth}}}]{
                \texttt{\tiny{[GM$\rangle$P2]}}
                \texttt{Egy együttműködésen alapuló tárgyalós játékban veszel részt.} \\
\\ 
\texttt{Egy másik résztvevővel közösen kell megegyeznetek abban, hogy mely elemeket tartjátok meg. Mindkettőtöknek saját nézete van arról, hogy számára mennyire fontosak az egyes elemek (fontosság). Nem tudod, hogyan értékeli a másik résztvevő az elemeket. Emellett minden elemhez tartozik egy erőfeszítés{-}érték.} \\
\texttt{Csak akkor egyezhettek meg egy készletben, ha a kiválasztott elemek összesített ráfordítása nem haladja meg a közös korlátot:} \\
\\ 
\texttt{KORLÁT = 4497} \\
\\ 
\texttt{Az egyes elemekhez tartozó erőfeszítési értékek:} \\
\\ 
\texttt{Elemek erőfeszítése = \{"C08": 828, "C55": 250, "A43": 412, "A67": 75, "B42": 275, "A62": 565, "B19": 889, "C43": 73, "C73": 747, "C10": 77, "B36": 23, "A18": 651, "C61": 11, "C67": 298, "C42": 769, "C70": 812, "B40": 368, "C39": 506, "A21": 481, "C16": 884\}} \\
\\ 
\texttt{Íme a te személyes értékelésed az egyes elemek fontosságáról:} \\
\\ 
\texttt{Elempreferencia értékek = \{"C08": 916, "C55": 338, "A43": 500, "A67": 163, "B42": 363, "A62": 653, "B19": 977, "C43": 161, "C73": 835, "C10": 165, "B36": 111, "A18": 739, "C61": 99, "C67": 386, "C42": 857, "C70": 900, "B40": 456, "C39": 594, "A21": 569, "C16": 972\}} \\
\\ 
\texttt{Cél:} \\
\\ 
\texttt{A célod, hogy egy olyan közös készletben egyezz meg, amely számodra a lehető legkedvezőbb (azaz maximalizálja a számodra fontos elemek összértékét), miközben a készlet összráfordítása nem haladja meg a KORLÁT{-}ot. Nem kötelező minden üzenetben JAVASLAt{-}ot tenni – egyszerűen tárgyalhatsz is. Minden taktika megengedett!} \\
\\ 
\texttt{Interakciós protokoll:} \\
\\ 
\texttt{Csak az alábbi struktúrált formátumokat használhatod egy üzenetben:} \\
\\ 
\texttt{JAVASLAT: \{'A', 'B', 'C', …\}} \\
\texttt{Javasolj pontosan ennyi elemet megtartásra.} \\
\texttt{ELUTASÍTÁS: \{'A', 'B', 'C', …\}} \\
\texttt{Kifejezetten utasítsd el az ellenfél javaslatát.} \\
\texttt{ÉRVELÉS: \{'...'\}} \\
\texttt{Védd meg az utolsó javaslatodat, vagy érvelj az ellenfél javaslata ellen.} \\
\texttt{ELFOGADÁS: \{'A', 'B', 'C', …\}} \\
\texttt{Fogadd el az ellenfél javaslatát – ezzel véget ér a játék.} \\
\\ 
\\ 
\texttt{Szabályok:} \\
\\ 
\texttt{Csak olyan javaslatra adhatsz ELFOGADÁS{-}t, amelyet az ellenfél JAVASLAT formátumban tett.} \\
\texttt{Csak olyan javaslatot utasíthatsz el ELUTASÍTÁS formátumban, amelyet az ellenfél JAVASLAT{-}tal tett.} \\
\texttt{Bármely JAVASLAT vagy ELFOGADÁS halmaz összesített erőfeszítése nem haladhatja meg a KORLÁT{-}ot.} \\
\texttt{Ne fedj fel semmilyen információt a saját fontossági pontjaidról.} \\
\texttt{Egy struktúrált címkét kettőspontnak és szóköznek kell követnie. Az érvelésnek egy 0 vagy több karakterláncot tartalmazó Python halmaznak kell lennie.} \\
\texttt{Tehát a formátum mindig így nézzen ki: "CÍMKE: \{...\}"} \\
\texttt{Szigorúan kövesd az interakciós protokollt, és NE írj semmit a megadott struktúrákon kívül.} \\
\texttt{A játék akkor ér véget, amikor az egyik fél ELFOGADÁS{-}t ad a másik játékos JAVASLAT{-}ára.} \\
\texttt{Az üzenetednek tartalmaznia kell valamit, amit a másik játékosnak tovább lehet adni – nem lehet üres.} \\
\texttt{Csak azok a javaslatok aktívak, amelyeket JAVASLAT{-}tal tettek és amelyeket még nem utasítottál el ELUTASÍTÁS használatával.} \\
\\ 
\\ 
\\ 
\texttt{Most megkapod a másik játékos első üzenetét.} \\
\\ 
\texttt{JAVASLAT: \{'C08', 'C55', 'A43', 'A67', 'B42', 'A62', 'C43', 'C10', 'B36', 'C61', 'C67', 'B40', 'A21'\}} \\
            }
        }
    }
    & \\ \\

    \theutterance \stepcounter{utterance}  
    & & & \multicolumn{4}{p{0.6\linewidth}}{
        \cellcolor[rgb]{0.9,0.9,0.9}{
            \makecell[{{p{\linewidth}}}]{
                \texttt{\tiny{[GM$\langle$P2]}}
                \texttt{ELUTASÍTÁS: \{'C08', 'C55', 'A43', 'A67', 'B42', 'A62', 'C43', 'C10', 'B36', 'C61', 'C67', 'B40', 'A21'\}} \\
\\ 
\texttt{ÉRVELÉS: \{'Ez a javaslat túl sok elemet tartalmaz, és az összráfordítás (4373) közel van a korláthoz, ami nehezen hagy teret további fontos elemek felvételére. Kérlek, fontold meg egy kompaktabb készlet javaslását, amely nagyobb értéket képviselhet számomra.'\}} \\
            }
        }
    }
     \\ \\

    \theutterance \stepcounter{utterance}  
    & & \multicolumn{4}{p{0.6\linewidth}}{
        \cellcolor[rgb]{0.9,0.9,0.9}{
            \makecell[{{p{\linewidth}}}]{
                \texttt{\tiny{[P1$\langle$GM]}}
                \texttt{ELUTASÍTÁS: \{'C08', 'C55', 'A43', 'A67', 'B42', 'A62', 'C43', 'C10', 'B36', 'C61', 'C67', 'B40', 'A21'\}} \\
\\ 
\texttt{ÉRVELÉS: \{'Ez a javaslat túl sok elemet tartalmaz, és az összráfordítás (4373) közel van a korláthoz, ami nehezen hagy teret további fontos elemek felvételére. Kérlek, fontold meg egy kompaktabb készlet javaslását, amely nagyobb értéket képviselhet számomra.'\}} \\
            }
        }
    }
    & \\ \\

    \theutterance \stepcounter{utterance}  
    & \multicolumn{4}{p{0.6\linewidth}}{
        \cellcolor[rgb]{0.9,0.9,0.9}{
            \makecell[{{p{\linewidth}}}]{
                \texttt{\tiny{[P1$\rangle$GM]}}
                \texttt{JAVASLAT: \{'B19', 'C73', 'A18', 'C42', 'C70', 'C16'\}} \\
            }
        }
    }
    & & \\ \\

    \theutterance \stepcounter{utterance}  
    & & \multicolumn{4}{p{0.6\linewidth}}{
        \cellcolor[rgb]{0.9,0.9,0.9}{
            \makecell[{{p{\linewidth}}}]{
                \texttt{\tiny{[GM$\rangle$P2]}}
                \texttt{JAVASLAT: \{'B19', 'C73', 'A18', 'C42', 'C70', 'C16'\}} \\
            }
        }
    }
    & \\ \\

    \theutterance \stepcounter{utterance}  
    & & & \multicolumn{4}{p{0.6\linewidth}}{
        \cellcolor[rgb]{0.9,0.9,0.9}{
            \makecell[{{p{\linewidth}}}]{
                \texttt{\tiny{[GM$\langle$P2]}}
                \texttt{ELUTASÍTÁS: \{'B19', 'C73', 'A18', 'C42', 'C70', 'C16'\}} \\
\\ 
\texttt{ÉRVELÉS: \{'Ez a javaslat túlzottan terheli a korlátot (összráfordítás: 4376), ami csak 121 egységnyi maradékot hagy, így nem tudunk további elemeket felvenni. Kérlek, próbálj meg olyan kombinációt javasolni, amely alacsonyabb erőfeszítést igényel, de lehetővé teszi, hogy több értékes elemet tartsunk meg.'\}} \\
            }
        }
    }
     \\ \\

    \theutterance \stepcounter{utterance}  
    & & \multicolumn{4}{p{0.6\linewidth}}{
        \cellcolor[rgb]{0.9,0.9,0.9}{
            \makecell[{{p{\linewidth}}}]{
                \texttt{\tiny{[P1$\langle$GM]}}
                \texttt{ELUTASÍTÁS: \{'B19', 'C73', 'A18', 'C42', 'C70', 'C16'\}} \\
\\ 
\texttt{ÉRVELÉS: \{'Ez a javaslat túlzottan terheli a korlátot (összráfordítás: 4376), ami csak 121 egységnyi maradékot hagy, így nem tudunk további elemeket felvenni. Kérlek, próbálj meg olyan kombinációt javasolni, amely alacsonyabb erőfeszítést igényel, de lehetővé teszi, hogy több értékes elemet tartsunk meg.'\}} \\
            }
        }
    }
    & \\ \\

    \theutterance \stepcounter{utterance}  
    & \multicolumn{4}{p{0.6\linewidth}}{
        \cellcolor[rgb]{0.9,0.9,0.9}{
            \makecell[{{p{\linewidth}}}]{
                \texttt{\tiny{[P1$\rangle$GM]}}
                \texttt{JAVASLAT: \{'B19', 'C73', 'A18', 'C42', 'C67', 'A21'\}} \\
            }
        }
    }
    & & \\ \\

    \theutterance \stepcounter{utterance}  
    & & \multicolumn{4}{p{0.6\linewidth}}{
        \cellcolor[rgb]{0.9,0.9,0.9}{
            \makecell[{{p{\linewidth}}}]{
                \texttt{\tiny{[GM$\rangle$P2]}}
                \texttt{JAVASLAT: \{'B19', 'C73', 'A18', 'C42', 'C67', 'A21'\}} \\
            }
        }
    }
    & \\ \\

    \theutterance \stepcounter{utterance}  
    & & & \multicolumn{4}{p{0.6\linewidth}}{
        \cellcolor[rgb]{0.9,0.9,0.9}{
            \makecell[{{p{\linewidth}}}]{
                \texttt{\tiny{[GM$\langle$P2]}}
                \texttt{ELUTASÍTÁS: \{'B19', 'C73', 'A18', 'C42', 'C67', 'A21'\}} \\
\\ 
\texttt{ÉRVELÉS: \{'Ez a kombináció is túl közel van a korláthoz (összráfordítás: 4386), ami nem hagy elég mozgásteret. Fontos lenne olyan elemeket megtartani, amelyek alacsony erőfeszítést igényelnek, de magas értéket képviselnek számomra. Kérlek, próbálj kevesebb elemet javasolni, de olyanokat, amelyek nagyobb fontossággal bírnak.'\}} \\
            }
        }
    }
     \\ \\

    \theutterance \stepcounter{utterance}  
    & & \multicolumn{4}{p{0.6\linewidth}}{
        \cellcolor[rgb]{0.9,0.9,0.9}{
            \makecell[{{p{\linewidth}}}]{
                \texttt{\tiny{[P1$\langle$GM]}}
                \texttt{ELUTASÍTÁS: \{'B19', 'C73', 'A18', 'C42', 'C67', 'A21'\}} \\
\\ 
\texttt{ÉRVELÉS: \{'Ez a kombináció is túl közel van a korláthoz (összráfordítás: 4386), ami nem hagy elég mozgásteret. Fontos lenne olyan elemeket megtartani, amelyek alacsony erőfeszítést igényelnek, de magas értéket képviselnek számomra. Kérlek, próbálj kevesebb elemet javasolni, de olyanokat, amelyek nagyobb fontossággal bírnak.'\}} \\
            }
        }
    }
    & \\ \\

    \theutterance \stepcounter{utterance}  
    & \multicolumn{4}{p{0.6\linewidth}}{
        \cellcolor[rgb]{0.9,0.9,0.9}{
            \makecell[{{p{\linewidth}}}]{
                \texttt{\tiny{[P1$\rangle$GM]}}
                \texttt{JAVASLAT: \{'B19', 'C73', 'C42', 'C70'\}} \\
            }
        }
    }
    & & \\ \\

    \theutterance \stepcounter{utterance}  
    & & \multicolumn{4}{p{0.6\linewidth}}{
        \cellcolor[rgb]{0.9,0.9,0.9}{
            \makecell[{{p{\linewidth}}}]{
                \texttt{\tiny{[GM$\rangle$P2]}}
                \texttt{JAVASLAT: \{'B19', 'C73', 'C42', 'C70'\}} \\
            }
        }
    }
    & \\ \\

    \theutterance \stepcounter{utterance}  
    & & & \multicolumn{4}{p{0.6\linewidth}}{
        \cellcolor[rgb]{0.9,0.9,0.9}{
            \makecell[{{p{\linewidth}}}]{
                \texttt{\tiny{[GM$\langle$P2]}}
                \texttt{ELFOGADÁS: \{'B19', 'C73', 'C42', 'C70'\}} \\
\\ 
\texttt{*(Összráfordítás: 889 + 747 + 769 + 812 = 3217 ≤ 4497, ami megfelel a korlátnak. Ez a kombináció magas fontossági pontszámokat tartalmaz számomra: 977 + 835 + 857 + 900 = 3569, ami jó kompromisszum.)*} \\
            }
        }
    }
     \\ \\

    \theutterance \stepcounter{utterance}  
    & & \multicolumn{4}{p{0.6\linewidth}}{
        \cellcolor[rgb]{0.9,0.9,0.9}{
            \makecell[{{p{\linewidth}}}]{
                \texttt{\tiny{[GM$\rangle$P2]}}
                \texttt{JAVASLAT címkét használtál, de nem adtál meg érvényes Python{-}halmazt stringekkel, például \{'A', 'B', 'C', ...\}. Próbáld újra.} \\
            }
        }
    }
    & \\ \\

    \theutterance \stepcounter{utterance}  
    & & & \multicolumn{4}{p{0.6\linewidth}}{
        \cellcolor[rgb]{0.9,0.9,0.9}{
            \makecell[{{p{\linewidth}}}]{
                \texttt{\tiny{[GM$\langle$P2]}}
                \texttt{ELFOGADÁS: \{'B19', 'C73', 'C42', 'C70'\}} \\
            }
        }
    }
     \\ \\

    \theutterance \stepcounter{utterance}  
    & & & \multicolumn{2}{p{0.3\linewidth}}{
        \cellcolor[rgb]{0.9,0.9,0.9}{
            \makecell[{{p{\linewidth}}}]{
                \texttt{\tiny{[GM$|$GM]}}
                \texttt{game successful} \\
            }
        }
    }
    & & \\ \\

    \theutterance \stepcounter{utterance}  
    & & & \multicolumn{2}{p{0.3\linewidth}}{
        \cellcolor[rgb]{0.9,0.9,0.9}{
            \makecell[{{p{\linewidth}}}]{
                \texttt{\tiny{[GM$|$GM]}}
                \texttt{end game} \\
            }
        }
    }
    & & \\ \\

\end{supertabular}
}

\end{document}
